% \iffalse meta-comment
% LTeX: language=de-DE
%
%  TUDCD-Script -- Corporate Design of Technische Universität Dresden
% ----------------------------------------------------------------------------
%
%  Copyright (C) TU Dresden <tudlatex@tu-dresden.de>, 2026
%
% ----------------------------------------------------------------------------
%
%  This work may be distributed and/or modified under the conditions of the
%  LaTeX Project Public License, either version 1.3c of this license or
%  any later version. The latest version of this license is in
%    http://www.latex-project.org/lppl.txt
%  and version 1.3c or later is part of all distributions of
%  LaTeX version 2008-05-04 or later.
%
%  This work has the LPPL maintenance status "maintained".
%
%  The current maintainer and author of this work is Yoshi Diepelt.
%
% ----------------------------------------------------------------------------
% \fi
%
% \iffalse ins:batch + dtx:driver
%<*ins>
\ifx\documentclass\undefined
  \input docstrip.tex
  \ifToplevel{\batchinput{tudcd.ins}}
\else
  \let\endbatchfile\relax
\fi
\endbatchfile
%</ins>
%<*dtx>
\ProvidesFile{tudcd-thesis.dtx}[2026/10/27]
\documentclass[english,ngerman]{tudcddoc}
\usepackage[T1]{fontenc}
\usepackage[ngerman=ngerman-x-latest]{hyphsubst}

\usepackage{babel}
\usepackage[babel]{microtype}
\RecordChanges
\begin{document} % Diese Dokumentation dokumentiert NUR diese Datei
  \title{\Large Dokumentation der Datei \texttt{\jobname.dtx} \\
  \normalsize Generiert durch \texttt{\$ enginetex \jobname.dtx}}
  \author{Vincent Voigtländer}
  \maketitle
  \tableofcontents

  \DocInput{tudcd-thesis.dtx}
\end{document}
%</dtx>
% \fi
% \selectlanguage{ngerman}
%
% \section{\texttt{tudcd-thesis}: Die Klasse für Abschlussarbeiten}
%
% \subsection{Vorbemerkung}
%
% Die Klasse \dcode{tudcdthesis} ist eine schlanke Hülle um die \KOMAScript-Klasse \dcode{scrbook}.
% Sie ersetzt die Klasse \dcode{tudscrreprt} aus dem \dpkg{tudscr}-Bündel und übernimmt
% deren inhaltliche Schnittstelle (Eingabefelder, Titelseite, Zusammenfassung und Erklärungen),
% nicht jedoch deren Layoutoptionen.
%
% Im Gegensatz zu den übrigen Klassen des \TUDCDScript{} werden
% \begin{itemize}
%   \item der Satzspiegel durch \dpkg{typearea} bestimmt (\doption{BCOR} und \doption{DIV} sind erlaubt),
%   \item der Fließtext in Computer Modern (Latin Modern) gesetzt und lediglich die Auszeichnungsschriften
%         (Überschriften, Titel, Kolumnentitel) in Noto Sans, sowie
%   \item die Seitenstile von \dpkg{scrlayer-scrpage} ohne Logo verwendet.
% \end{itemize}
%
% Aus den gemeinsamen Modulen werden lediglich \texttt{tudcd-common} (Optionen, Logo),
% \texttt{tudcd-documents} (Identifikation, Basisklasse) und \texttt{tudcd-colors} (Farben) genutzt.
%
%    \begin{macrocode}
%<@@=tudcd>
%    \end{macrocode}
%
% \subsection{Optionen}
%
% \subsubsection{Ausgabemedium}
%
% \begin{codedescribe}[option,new=0.8.0]{output}
% \begin{codesyntax}
% \tsobj[key]{output}=\tsobj[value]{digital,print}
% \end{codesyntax}
% Die globale Option \doption{output} legt fest, ob die Arbeit für die Ansicht am Bildschirm
% (\dcode{digital}, Voreinstellung) oder für den Druck (\dcode{print}) gesetzt wird.
%
% Für \dcode{digital} wird einseitig (\doption{oneside}) gesetzt und alle Elemente des Corporate Designs
% erscheinen in Brilliantblau.
%
% Für \dcode{print} werden die \KOMAScript-Optionen \doption{twoside}, \doption{open}\dcode{=right},
% \doption{cleardoublepage}\dcode{=empty} und \doption{BCOR}\dcode{=10mm} voreingestellt.
% Zudem werden Überschriften, Titel, Kolumnentitel, Verweise und das Logo schwarz gesetzt,
% um farbige Druckseiten zu vermeiden.
%
% Die Voreinstellungen werden nur dann gesetzt, wenn die jeweilige Option nicht explizit
% in \dmacro{\documentclass} angegeben worden ist.
% \end{codedescribe}
%
% Die Wahl wird in einer Zeichenkette gespeichert.
% Weiterhin wird eine Liste der durch Nutzende explizit gesetzten \KOMAScript-Optionen benötigt,
% welche in \texttt{tudcd-documents} befüllt wird.
%    \begin{macrocode}
%<*prelim-declaration>
\str_new:N \l_@@_thesis_output_str
\clist_new:N \g_@@_thesis_user_koma_clist
\bool_new:N \l_@@_thesis_bodyfont_noto_bool
%</prelim-declaration>
%    \end{macrocode}
%
% Wird \dcode{print} gewählt, werden die Hervorhebungsfarbe und die Farbe des Logos auf Schwarz gesetzt.
% Werden \doption{color/highlight-color} oder \doption{logo/color} später gesetzt, so überschreiben diese die Voreinstellung.
%    \begin{macrocode}
%<*option&peri>
  output .choices:nn = { digital, print }
    {
      \str_set:NV \l_@@_thesis_output_str \l_keys_choice_tl
      \str_if_eq:VnT \l_keys_choice_tl { print }
        {
          \tl_set:Nn \l_@@_highlight_color_tl { Schwarz }
          \str_set:Nn \l_@@_logocolor_str { black }
        }
    },
  output .initial:n = digital,
  output .usage:n = load,
%</option&peri>
%    \end{macrocode}
%
% \subsubsection{Schrift des Fließtextes}
%
% \begin{codedescribe}[option,new=0.8.0]{bodyfont}
% \begin{codesyntax}
% \tsobj[key]{bodyfont}=\tsobj[value]{cm,noto}
% \end{codesyntax}
% Der Fließtext wird per Voreinstellung (\dcode{cm}) in Computer Modern gesetzt, da diese Schrift
% für lange Texte gut lesbar ist. Mit \dcode{noto} wird auch der Fließtext in Noto Sans gesetzt.
% Der Mathematiksatz bleibt in beiden Fällen unverändert.
% \end{codedescribe}
%    \begin{macrocode}
%<*option&peri>
  bodyfont .choices:nn = { cm, noto }
    {
      \str_if_eq:VnTF \l_keys_choice_tl { noto }
        { \bool_set_true:N \l_@@_thesis_bodyfont_noto_bool }
        { \bool_set_false:N \l_@@_thesis_bodyfont_noto_bool }
    },
  bodyfont .initial:n = cm,
  bodyfont .usage:n = preamble,
%</option&peri>
%    \end{macrocode}
%
% \subsubsection{Zusammenfassung und Erklärungen}
%
% \begin{codedescribe}[option,new=0.8.0]{abstract,declaration}
% \begin{codesyntax}
% \tsobj[key]{abstract}=\tsobj[value]{Werteliste}
%
% \tsobj[key]{declaration}=\tsobj[value]{Werteliste}
% \end{codesyntax}
% Die beiden Optionen steuern das Aussehen der Umgebungen \denv{abstract} und \denv{declarations}.
% Sie verstehen dieselben Werte wie die gleichnamigen Optionen von \dpkg{tudscr}, inklusive aller Synonyme:
% \begin{description}
%   \item[\dcode{false}] keine Überschrift,
%   \item[\dcode{true}] zentrierte Überschrift (Voreinstellung),
%   \item[\dcode{section}] Überschrift mit \dmacro{\section*},
%   \item[\dcode{chapter}] Überschrift mit \dmacro{\chapter*},
%   \item[\dcode{single}, \dcode{multiple}] jeder Teil auf eigener Seite oder alle Teile gemeinsam,
%   \item[\dcode{fil}, \dcode{nofil}] vertikale Zentrierung auf der Seite,
%   \item[\dcode{markboth}, \dcode{nomarkboth}] Setzen der Kolumnentitel,
%   \item[\dcode{nottotoc}, \dcode{totoc}, \dcode{totocleveldown}, \dcode{tocmultiple}] Einträge im Inhaltsverzeichnis.
% \end{description}
% \end{codedescribe}
%
% Für beide Optionen wird ein gemeinsamer Zustand angelegt.
%    \begin{macrocode}
%<*prelim-declaration>
\clist_map_inline:nn { abstract, declaration }
  {
    \int_new:c { l_@@_thesis_#1_level_int }
    \int_new:c { l_@@_thesis_#1_toc_int }
    \bool_new:c { l_@@_thesis_#1_toc_locked_bool }
    \bool_new:c { l_@@_thesis_#1_multi_bool }
    \bool_new:c { l_@@_thesis_#1_fil_bool }
    \bool_new:c { l_@@_thesis_#1_markboth_bool }
    \int_set:cn { l_@@_thesis_#1_level_int } { 1 }
    \bool_set_true:c { l_@@_thesis_#1_multi_bool }
    \bool_set_true:c { l_@@_thesis_#1_fil_bool }
  }
%    \end{macrocode}
% Die Synonyme der Werte werden auf einen kanonischen Wert abgebildet.
%    \begin{macrocode}
\prop_new:N \g_@@_thesis_fm_alias_prop
\cs_generate_variant:Nn \prop_get:NnNTF { Ne }
\cs_set_protected:Npn \@@_thesis_tmp:nn #1#2
  {
    \clist_map_inline:nn {#2}
      { \prop_gput:Nnn \g_@@_thesis_fm_alias_prop {##1} {#1} }
  }
\@@_thesis_tmp:nn { false } { false, off, no }
\@@_thesis_tmp:nn { true } { true, on, yes }
\@@_thesis_tmp:nn { section } { section, sect, sec, addsec }
\@@_thesis_tmp:nn { chapter } { chapter, chap, addchap, heading, new, standard }
\@@_thesis_tmp:nn { single } { single, one, simple }
\@@_thesis_tmp:nn { multiple } { multiple, multi, all, aggregate, two, both, double }
\@@_thesis_tmp:nn { nofil } { nofil, nofill, novfil, novfill }
\@@_thesis_tmp:nn { fil } { fil, fill, vfil, vfill }
\@@_thesis_tmp:nn { markboth } { markboth, mark }
\@@_thesis_tmp:nn { nomarkboth } { nomarkboth, nomark }
\@@_thesis_tmp:nn { nottotoc } { nottotoc, notoc }
\@@_thesis_tmp:nn { totoc } { totoc, toc, totoclevelup, toclevelup, levelup }
\@@_thesis_tmp:nn { totocleveldown } { totocleveldown, tocleveldown, leveldown }
\@@_thesis_tmp:nn { tocmultiple }
  {
    tocmultiple, tocmulti, tocall, tocaggregate,
    totocmultiple, totocmulti, totocall, totocaggregate
  }
%    \end{macrocode}
% Die Ebene der Überschrift gibt den Eintrag ins Inhaltsverzeichnis vor, sofern dieser nicht explizit gesetzt worden ist.
% Die Zahlenwerte entsprechen denen von \dpkg{tudscr}: Ebene $0$ bis $3$ für \dcode{false} bis \dcode{chapter},
% Inhaltsverzeichnis $0$ (kein Eintrag), $1$ (oberste Ebene), $2$ (eine Ebene tiefer) sowie $3$ und $4$ für
% einen gemeinsamen Eintrag aller Teile.
%    \begin{macrocode}
\msg_new:nnnn { tudcd } { thesis/fm-value }
  { Unknown~value~'#2'~for~option~'#1'. }
  { See~the~tudcdscr~handbook~for~the~values~of~the~'#1'~option. }
\cs_new_protected:Npn \@@_thesis_fm_toc_preset:nn #1#2
  {
    \bool_if:cF { l_@@_thesis_#1_toc_locked_bool }
      { \int_set:cn { l_@@_thesis_#1_toc_int } {#2} }
  }
\cs_new_protected:Npn \@@_thesis_fm_toc_set:nn #1#2
  {
    \int_set:cn { l_@@_thesis_#1_toc_int } {#2}
    \bool_set_true:c { l_@@_thesis_#1_toc_locked_bool }
  }
\cs_new_protected:Npn \@@_thesis_fm_set:nn #1#2
  {
    \clist_map_inline:nn {#2}
      {
        \prop_get:NeNTF \g_@@_thesis_fm_alias_prop { \tl_to_str:n {##1} } \l_@@_thesis_tmp_tl
          { \@@_thesis_fm_apply:nV {#1} \l_@@_thesis_tmp_tl }
          { \msg_error:nnnn { tudcd } { thesis/fm-value } {#1} {##1} }
      }
  }
\tl_new:N \l_@@_thesis_tmp_tl
\cs_new_protected:Npn \@@_thesis_fm_apply:nn #1#2
  {
    \str_case:nn {#2}
      {
        { false }
          {
            \int_set:cn { l_@@_thesis_#1_level_int } { 0 }
            \@@_thesis_fm_toc_preset:nn {#1} { 0 }
          }
        { true }
          {
            \int_set:cn { l_@@_thesis_#1_level_int } { 1 }
            \@@_thesis_fm_toc_preset:nn {#1} { 0 }
          }
        { section }
          {
            \int_set:cn { l_@@_thesis_#1_level_int } { 2 }
            \@@_thesis_fm_toc_preset:nn {#1} { 1 }
          }
        { chapter }
          {
            \int_set:cn { l_@@_thesis_#1_level_int } { 3 }
            \@@_thesis_fm_toc_preset:nn {#1} { 1 }
          }
        { single } { \bool_set_false:c { l_@@_thesis_#1_multi_bool } }
        { multiple } { \bool_set_true:c { l_@@_thesis_#1_multi_bool } }
        { nofil } { \bool_set_false:c { l_@@_thesis_#1_fil_bool } }
        { fil } { \bool_set_true:c { l_@@_thesis_#1_fil_bool } }
        { markboth } { \bool_set_true:c { l_@@_thesis_#1_markboth_bool } }
        { nomarkboth } { \bool_set_false:c { l_@@_thesis_#1_markboth_bool } }
        { nottotoc } { \@@_thesis_fm_toc_set:nn {#1} { 0 } }
        { totoc }
          {
            \@@_thesis_fm_toc_set:nn {#1}
              { \int_compare:nNnTF { \int_use:c { l_@@_thesis_#1_toc_int } } < { 3 } { 1 } { 3 } }
          }
        { totocleveldown }
          {
            \@@_thesis_fm_toc_set:nn {#1}
              { \int_compare:nNnTF { \int_use:c { l_@@_thesis_#1_toc_int } } < { 3 } { 2 } { 4 } }
          }
        { tocmultiple }
          {
            \int_compare:nNnT { \int_use:c { l_@@_thesis_#1_toc_int } } = { 0 }
              { \int_set:cn { l_@@_thesis_#1_toc_int } { 1 } }
            \int_compare:nNnT { \int_use:c { l_@@_thesis_#1_toc_int } } < { 3 }
              { \int_add:cn { l_@@_thesis_#1_toc_int } { 2 } }
            \bool_set_true:c { l_@@_thesis_#1_toc_locked_bool }
          }
      }
  }
\cs_generate_variant:Nn \@@_thesis_fm_apply:nn { nV }
%</prelim-declaration>
%    \end{macrocode}
%
%    \begin{macrocode}
%<*option&peri>
  abstract .code:n = { \@@_thesis_fm_set:nn { abstract } {#1} },
  abstract .default:n = true,
  declaration .code:n = { \@@_thesis_fm_set:nn { declaration } {#1} },
  declaration .default:n = true,
%</option&peri>
%    \end{macrocode}
%
% \subsubsection{Sprache des Logos}
%
% Die Sprache des Logos folgt per Voreinstellung der Hauptsprache des Dokuments.
% Dafür wird vor dem Verarbeiten der Optionen \doption{logo/language}\dcode{=auto} gesetzt.
% Wählen Nutzende explizit eine Sprache, so wird die automatische Erkennung wieder entfernt.
%    \begin{macrocode}
%<*optionloading&pre>
\str_clear:N \l_@@_logolang_str
\keys_set:nn { tudcd } { logo/language = auto }
%</optionloading&pre>
%<*optionloading&post>
\str_if_empty:NF \l_@@_logolang_str
  {
    \hook_gremove_code:nn { package/babel/after } { tudcd-language-detection-babel }
    \hook_gremove_code:nn { package/polyglossia/after } { tudcd-language-detection-polyglossia }
  }
\hook_gput_code:nnn { begindocument/before } { tudcd/thesis/logolanguage }
  {
    \str_if_empty:NTF \l_@@_logolang_str
      { \str_set:Nn \l_@@_logolang_str { de } }
      {
        \str_case:VnF \l_@@_logolang_str { { de } { } { en } { } }
          { \str_set:Nn \l_@@_logolang_str { en } }
      }
  }
%</optionloading&post>
%    \end{macrocode}
%
% \subsubsection{Voreinstellungen für \KOMAScript}
%
% Die Voreinstellungen der Option \doption{output} werden vor dem Laden von \dcode{scrbook} weitergereicht.
% Da \KOMAScript{} die globalen Optionen aus \dmacro{\documentclass} vor den weitergereichten Optionen verarbeitet,
% wird eine Voreinstellung nur dann weitergereicht, wenn keiner der zugehörigen Schlüssel explizit gesetzt worden ist.
%    \begin{macrocode}
%<*optionloading&post>
\cs_new_protected:Npn \@@_thesis_koma_default:nn #1#2
  {
    \bool_set_false:N \l_tmpa_bool
    \clist_map_inline:Nn \g_@@_thesis_user_koma_clist
      {
        \clist_map_inline:nn {#1}
          { \str_if_eq:nnT {##1} {####1} { \bool_set_true:N \l_tmpa_bool } }
      }
    \bool_if:NF \l_tmpa_bool { \PassOptionsToClass {#2} { scrbook } }
  }
\str_if_eq:VnTF \l_@@_thesis_output_str { print }
  {
    \@@_thesis_koma_default:nn { oneside, twoside } { twoside }
    \@@_thesis_koma_default:nn { open, openright, openany, openleft } { open=right }
    \@@_thesis_koma_default:nn { cleardoublepage } { cleardoublepage=empty }
    \@@_thesis_koma_default:nn { BCOR } { BCOR=10mm }
  }
  { \@@_thesis_koma_default:nn { oneside, twoside } { oneside } }
%</optionloading&post>
%    \end{macrocode}
%
% \subsection{Schriften}
%
% Der Fließtext bleibt in Computer Modern. Unter pdf\TeX{} wird dafür \dpkg{lmodern} mit der Kodierung T1 geladen,
% die Auszeichnungsschrift Noto Sans stellt das Paket \dpkg{noto-sans} bereit.
% Unter Lua\TeX{} und Xe\TeX{} ist Latin Modern bereits die Voreinstellung von \dpkg{fontspec},
% Noto Sans wird direkt über die Dateinamen eingebunden.
% Das Paket \dpkg{noto-sans} wird dort bewusst nicht genutzt, da es globale Schrifteigenschaften
% (\dmacro{\defaultfontfeatures}) setzt, die von Nutzenden geladene Schriften beeinflussen würden.
%
% Das Paket \dpkg{iftex} wird wie bei \dpkg{tudscr} immer geladen, da viele Vorlagen \dmacro{\iftutex} verwenden.
%    \begin{macrocode}
%<*body&peri>
\RequirePackage{iftex}
\sys_if_engine_opentype:TF
  {
    \RequirePackage{fontspec}
    \setsansfont { NotoSans }
      [
        Extension = .ttf ,
        UprightFont = *-Regular ,
        ItalicFont = *-Italic ,
        BoldFont = *-SemiBold ,
        BoldItalicFont = *-SemiBoldItalic
      ]
  }
  {
    \RequirePackage[T1]{fontenc}
    \RequirePackage{lmodern}
    \RequirePackage[semibold]{noto-sans}
  }
\hook_gput_code:nnn { begindocument/before } { tudcd/thesis/bodyfont }
  {
    \bool_if:NT \l_@@_thesis_bodyfont_noto_bool
      {
        \renewcommand*{\familydefault}{\sfdefault}
        \normalfont
      }
  }
%</body&peri>
%    \end{macrocode}
%
% Im Vorspann (\cs{frontmatter}) werden die Seiten mit großen römischen Ziffern nummeriert.
%    \begin{macrocode}
%<*body&peri>
\hook_gput_code:nnn { begindocument/before } { tudcd/thesis/romannumbers }
  {
    \cs_if_exist:NT \frontmatter
      {
        \hook_gput_code:nnn { cmd/frontmatter/after } { tudcd/thesis/romannumbers }
          { \pagenumbering { Roman } }
      }
  }
%</body&peri>
%    \end{macrocode}
%
% \subsection{Gestaltungselemente}
%
% Alle Gliederungsüberschriften und, da \dcode{scrbook} diese von \dcode{disposition} ableitet,
% auch die Kapiteleinträge im Inhaltsverzeichnis werden in Noto Sans und der Hervorhebungsfarbe gesetzt.
%    \begin{macrocode}
%<*body&peri>
\setkomafont{disposition}{\normalcolor\sffamily\bfseries\@@_select_highlight_color:}
%    \end{macrocode}
%
% Für die Titelseite werden die Elemente von \KOMAScript{} angepasst und zwei neue Elemente angelegt:
% \dcode{thesis} für die Art der Arbeit und den akademischen Grad sowie \dcode{titleinfo} für alle
% weiteren Angaben.
%    \begin{macrocode}
\setkomafont{title}{\usekomafont{disposition}\huge}
\setkomafont{subtitle}{\usekomafont{title}\Large\mdseries}
\setkomafont{subject}{\normalcolor\sffamily\Large}
\setkomafont{author}{\normalcolor\sffamily\Large}
\setkomafont{date}{\normalcolor\sffamily}
\setkomafont{publishers}{\normalcolor\sffamily}
\newkomafont{thesis}{\usekomafont{disposition}\Large}
\newkomafont{titleinfo}{\normalcolor\sffamily}
%    \end{macrocode}
%
% \subsection{Seitenstile}
%
% Es werden die Seitenstile von \dpkg{scrlayer-scrpage} verwendet.
% Kopf- und Fußzeile erhalten Noto Sans und die Hervorhebungsfarbe.
% Für Dokumente, welche für \dcode{tudscrreprt} geschrieben wurden, werden die dortigen Seitenstile
% als Aliase bereitgestellt.
%    \begin{macrocode}
\RequirePackage{scrlayer-scrpage}
\setkomafont{pageheadfoot}{\normalcolor\sffamily\@@_select_highlight_color:}
\automark[section]{chapter}
\pagestyle{scrheadings}
\DeclarePageStyleAlias{tudheadings}{scrheadings}
\DeclarePageStyleAlias{plain.tudheadings}{plain.scrheadings}
\DeclarePageStyleAlias{empty.tudheadings}{empty}
%    \end{macrocode}
%
% Wird \dpkg{hyperref} geladen, erhalten Verweise die Hervorhebungsfarbe, im Druck also Schwarz.
% Eigene Einstellungen mittels \dmacro{\hypersetup} haben Vorrang.
%    \begin{macrocode}
\hook_gput_code:nnn { package/hyperref/after } { tudcd/thesis }
  { \hypersetup { colorlinks , allcolors = highlightcolor } }
%</body&peri>
%    \end{macrocode}
%
% \subsection{Optionen nach dem Laden der Klasse}
%
% \begin{codedescribe}[macro,new=0.8.0]{\TUDoptions,\TUDoption}
% \begin{codesyntax}
% \tsmacro{\TUDoptions}{Schlüssel-Wert-Liste}
%
% \tsmacro{\TUDoption}{Schlüssel,Wert}
% \end{codesyntax}
% Zur Kompatibilität mit \dpkg{tudscr} existieren beide Befehle als Alternativen zu \dmacro{\TUDCDsetup}.
% Unbekannte Schlüssel werden nach dem Laden der Klasse an \dmacro{\KOMAoptions} weitergereicht.
% Kennt auch \KOMAScript{} den Schlüssel nicht (bspw.\@ die Layoutoptionen von \dpkg{tudscr}),
% wird er mit einer Warnung ignoriert.
% \end{codedescribe}
%
% Ob \KOMAScript{} einen Schlüssel kennt, wird anhand der Makros \dcode{KV@<Mitglied>@<Schlüssel>}
% aller Mitglieder der Familie \dcode{KOMA} geprüft, wie es auch \dpkg{scrbase} selbst tut.
%    \begin{macrocode}
%<*body&peri>
\NewDocumentCommand \TUDoptions { m } { \keys_set:nn { tudcd } {#1} }
\NewDocumentCommand \TUDoption { m m } { \keys_set:nn { tudcd } { #1 = {#2} } }
\msg_new:nnnn { tudcd } { thesis/unknown-option }
  { Unknown~option~'#1'~ignored. }
  {
    '#1'~is~neither~an~option~of~tudcdthesis~nor~of~KOMA-Script.~
    The~layout~options~of~tudscr~are~not~supported~by~tudcdthesis.
  }
\prg_new_protected_conditional:Npnn \@@_thesis_if_koma_key:n #1 { TF }
  {
    \bool_set_false:N \l_tmpa_bool
    \clist_map_inline:en { KOMA, \cs_if_exist_use:c { KOMA@key@familylist } }
      { \cs_if_exist:cT { KV@##1@#1 } { \bool_set_true:N \l_tmpa_bool } }
    \bool_if:NTF \l_tmpa_bool { \prg_return_true: } { \prg_return_false: }
  }
\cs_generate_variant:Nn \clist_map_inline:nn { en }
\keys_define:nn { tudcd }
  {
    unknown .code:n =
      {
        \exp_args:NV \@@_thesis_if_koma_key:nTF \l_keys_key_str
          {
            \tl_if_novalue:nTF {#1}
              { \exp_args:Ne \KOMAoptions { \l_keys_key_str } }
              { \exp_args:Ne \KOMAoptions { \l_keys_key_str = \exp_not:n {#1} } }
          }
          { \msg_warning:nnV { tudcd } { thesis/unknown-option } \l_keys_key_str }
      }
  }
%</body&peri>
%    \end{macrocode}
%
% \subsection{Eingabefelder}
%
% Die Felder werden wie bei \dpkg{tudscr} in Makros der Form \dcode{\textbackslash @<Feld>} gespeichert,
% damit eigene Texte (beispielsweise mittels \dmacro{\renewcaptionname}) über \dmacro{\getfield} darauf zugreifen können.
%
% \begin{codedescribe}[macro,new=0.8.0]{\faculty,\department,\institute,\chair}
% \begin{codesyntax}
% \tsmacro{\faculty}[ignoriert]{Fakultät}
% \end{codesyntax}
% Die Organisationseinheiten werden auf der Titelseite unter dem Logo ausgegeben.
% Das optionale Argument wird aus Kompatibilitätsgründen akzeptiert und ignoriert.
% \end{codedescribe}
%    \begin{macrocode}
%<*body&peri>
\cs_new_protected:Npn \@@_thesis_new_field:nn #1#2
  {
    \tl_gclear_new:c { @#1 }
    \tl_gset:cn { @#1 } {#2}
  }
\clist_map_inline:nn { faculty, department, institute, chair }
  {
    \@@_thesis_new_field:nn {#1} { }
    \exp_args:Nc \NewDocumentCommand { #1 } { o m }
      { \tl_gset:cn { @#1 } {##2} }
  }
%    \end{macrocode}
%
% \begin{codedescribe}[macro,new=0.8.0]{\supervisor,\professor,\referee,\advisor,\company,\defensedate,\place,\confirmationclosing}
% \begin{codesyntax}
% \tsmacro{\supervisor}{Name \textbackslash and Name}
% \end{codesyntax}
% Betreuende, Hochschullehrende, Gutachtende und Fachreferenten werden mit \dmacro{\and} getrennt.
% \dmacro{\company} wird im Sperrvermerk, \dmacro{\place} (Voreinstellung: Dresden) und
% \dmacro{\confirmationclosing} in der Selbstständigkeitserklärung verwendet.
% \end{codedescribe}
%    \begin{macrocode}
\clist_map_inline:nn
  { supervisor, professor, referee, advisor, company, defensedate, place }
  {
    \@@_thesis_new_field:nn {#1} { }
    \exp_args:Nc \NewDocumentCommand { #1 } { m }
      { \tl_gset:cn { @#1 } {##1} }
  }
\tl_gset:Nn \@place { Dresden }
\@@_thesis_new_field:nn { supporter } { }
\NewDocumentCommand \confirmationclosing { m }
  { \tl_gset:Nn \@confirmationclosing {#1} }
%    \end{macrocode}
%
% \begin{codedescribe}[macro,new=0.8.0]{\thesis,\subject}
% \begin{codesyntax}
% \tsmacro{\thesis}{Art der Arbeit}
%
% \tsmacro{\subject}{Art der Arbeit}
% \end{codesyntax}
% Die Art der Arbeit kann als Freitext oder als eines der Schlüsselwörter
% \dcode{habil}, \dcode{diss} (\dcode{doctoral}, \dcode{phd}), \dcode{diploma}, \dcode{master},
% \dcode{bachelor}, \dcode{student}, \dcode{evidence}, \dcode{project}, \dcode{seminar},
% \dcode{term}, \dcode{research}, \dcode{log}, \dcode{report} und \dcode{internship}
% angegeben werden. Schlüsselwörter werden durch die sprachabhängige Bezeichnung ersetzt.
% Wird \dmacro{\subject} verwendet und ist noch keine Art der Arbeit gesetzt, so wird diese übernommen.
% \end{codedescribe}
%    \begin{macrocode}
\@@_thesis_new_field:nn { thesis } { }
\prop_const_from_keyval:Nn \c_@@_thesis_type_prop
  {
    habil = \habilitationname ,
    diss = \dissertationname ,
    doctoral = \dissertationname ,
    phd = \dissertationname ,
    diploma = \diplomathesisname ,
    master = \masterthesisname ,
    bachelor = \bachelorthesisname ,
    student = \studentthesisname ,
    evidence = \studentresearchname ,
    project = \projectpapername ,
    seminar = \seminarpapername ,
    term = \termpapername ,
    research = \researchname ,
    log = \logname ,
    report = \reportname ,
    internship = \internshipname
  }
\cs_new_protected:Npn \@@_thesis_type_set:Nn #1#2
  {
    \prop_get:NeNTF \c_@@_thesis_type_prop { \tl_trim_spaces:n {#2} } \l_@@_thesis_tmp_tl
      { \tl_gset_eq:NN #1 \l_@@_thesis_tmp_tl }
      { \tl_gset:Nn #1 {#2} }
  }
\NewDocumentCommand \thesis { m } { \@@_thesis_type_set:Nn \@thesis {#1} }
\RenewDocumentCommand \subject { m }
  {
    \@@_thesis_type_set:Nn \@subject {#1}
    \tl_if_empty:NT \@thesis { \tl_gset_eq:NN \@thesis \@subject }
  }
%    \end{macrocode}
%
% \begin{codedescribe}[macro,new=0.8.0]{\graduation}
% \begin{codesyntax}
% \tsmacro{\graduation}[Abkürzung]{Akademischer Grad}
% \end{codesyntax}
% \end{codedescribe}
%    \begin{macrocode}
\@@_thesis_new_field:nn { graduation } { }
\@@_thesis_new_field:nn { graduationabbreviation } { }
\NewDocumentCommand \graduation { o m }
  {
    \tl_gset:Nn \@graduation {#2}
    \IfValueTF {#1}
      { \tl_gset:Nn \@graduationabbreviation {#1} }
      { \tl_gclear:N \@graduationabbreviation }
  }
%    \end{macrocode}
%
% Der Titel wird zusätzlich in \dcode{\textbackslash tudcd@thesis@title} gespeichert, damit die Texte der Erklärungen
% prüfen können, ob ein Titel gesetzt worden ist.
%    \begin{macrocode}
\tl_new:N \tudcd@thesis@title
\RenewDocumentCommand \title { m }
  {
    \tl_gset:Nn \@title {#1}
    \tl_gset:Nn \tudcd@thesis@title {#1}
  }
%    \end{macrocode}
%
% \begin{codedescribe}[macro,new=0.8.0]{\getfield}
% \begin{codesyntax}
% \tsmacro{\getfield}[Befehl]{Feld}
% \end{codesyntax}
% Gibt den Inhalt eines Feldes aus, beispielsweise \dcode{\textbackslash getfield\{title\}}.
% Ist ein \tsobj[oarg]{Befehl} angegeben, wird diesem der Inhalt als Argument übergeben.
% \end{codedescribe}
%    \begin{macrocode}
\msg_new:nnn { tudcd } { thesis/unknown-field }
  { The~field~'#1'~is~unknown~and~cannot~be~printed. }
\NewExpandableDocumentCommand \getfield { o m }
  {
    \cs_if_exist:cTF { @#2 }
      { \IfValueTF {#1} { #1 { \use:c { @#2 } } } { \use:c { @#2 } } }
      { \msg_expandable_error:nnn { tudcd } { thesis/unknown-field } {#2} }
  }
%</body&peri>
%    \end{macrocode}
%
% \subsubsection{Angaben zu den Autoren}
%
% \begin{codedescribe}[macro,new=0.8.0]{\matriculationnumber,\matriculationyear,\dateofbirth,\placeofbirth,\course,\discipline,\emailaddress,\authormore}
% \begin{codesyntax}
% \tsmacro{\matriculationnumber}{Matrikelnummer}
%
% \tsmacro{\emailaddress}[ignoriert]{E-Mail-Adresse}
% \end{codesyntax}
% Die Angaben können außerhalb von \dmacro{\author} für alle Autoren oder innerhalb von \dmacro{\author}
% direkt nach dem jeweiligen Namen für einzelne Autoren gesetzt werden.
% Bei der gesternten Variante von \dmacro{\emailaddress} wird die Adresse nicht verlinkt.
% \end{codedescribe}
%
% Außerhalb von \dmacro{\author} speichern die Befehle den Wert global.
%    \begin{macrocode}
%<*body&peri>
\clist_const:Nn \c_@@_thesis_authorfields_clist
  {
    matriculationnumber, matriculationyear, dateofbirth, placeofbirth,
    course, discipline, emailaddress, authormore
  }
\clist_map_inline:Nn \c_@@_thesis_authorfields_clist
  { \@@_thesis_new_field:nn {#1} { } }
\bool_new:N \g_@@_thesis_emailaddress_nolink_bool
\clist_map_inline:nn
  {
    matriculationnumber, matriculationyear, dateofbirth, placeofbirth,
    course, discipline, authormore
  }
  {
    \exp_args:Nc \NewDocumentCommand { #1 } { m }
      { \tl_gset:cn { @#1 } {##1} }
  }
\NewDocumentCommand \emailaddress { s o m }
  {
    \tl_gset:Nn \@emailaddress {#3}
    \bool_gset:Nn \g_@@_thesis_emailaddress_nolink_bool {#1}
  }
%    \end{macrocode}
%
% Innerhalb von \dmacro{\author} wird die Liste beim Aufruf an \dmacro{\and} getrennt.
% Für jeden Autor werden die Angaben zweimal ausgewertet:
% Einmal in einer Box, in welcher die Befehle ihren Wert in einer Eigenschaftsliste je Autor ablegen,
% und einmal expandierend, wobei die Befehle samt Argument entfernt werden und nur der Name übrig bleibt.
% \dcode{\textbackslash @author} enthält anschließend nur die Namen, getrennt durch \dmacro{\and}.
%    \begin{macrocode}
\seq_new:N \g_@@_thesis_author_seq
\int_new:N \g_@@_thesis_author_int
\seq_new:N \l_@@_thesis_tmp_seq
\box_new:N \l_@@_thesis_tmp_box
\cs_new_protected:Npn \@@_thesis_author_collect:
  {
    \clist_map_inline:Nn \c_@@_thesis_authorfields_clist
      {
        \cs_set_protected:cpn {##1} ####1
          {
            \prop_gput:cnn
              { g_@@_thesis_author_ \int_use:N \g_@@_thesis_author_int _prop }
              {##1} {####1}
          }
      }
    \DeclareDocumentCommand \emailaddress { s o m }
      {
        \prop_gput:cnn
          { g_@@_thesis_author_ \int_use:N \g_@@_thesis_author_int _prop }
          { emailaddress } {##3}
        \IfBooleanT {##1}
          {
            \prop_gput:cnn
              { g_@@_thesis_author_ \int_use:N \g_@@_thesis_author_int _prop }
              { emailaddressnolink } { }
          }
      }
  }
\cs_new_protected:Npn \@@_thesis_author_strip:
  {
    \clist_map_inline:Nn \c_@@_thesis_authorfields_clist
      { \cs_set:cpn {##1} ####1 { } }
    \DeclareExpandableDocumentCommand \emailaddress { s o m } { }
  }
\cs_new_protected:Npn \@@_thesis_author_parse:n #1
  {
    \seq_gclear:N \g_@@_thesis_author_seq
    \int_gzero:N \g_@@_thesis_author_int
    \seq_set_split:Nnn \l_@@_thesis_tmp_seq { \and } {#1}
    \seq_map_inline:Nn \l_@@_thesis_tmp_seq
      {
        \int_gincr:N \g_@@_thesis_author_int
        \prop_gclear_new:c
          { g_@@_thesis_author_ \int_use:N \g_@@_thesis_author_int _prop }
        \group_begin:
          \@@_thesis_author_collect:
          \cs_set_eq:NN \thanks \use_none:n
          \hbox_set:Nn \l_@@_thesis_tmp_box {##1}
        \group_end:
        \group_begin:
          \@@_thesis_author_strip:
          \protected@edef \l_@@_thesis_tmp_tl {##1}
          \tl_set:Ne \l_@@_thesis_tmp_tl { \tl_trim_spaces:V \l_@@_thesis_tmp_tl }
          \seq_gput_right:NV \g_@@_thesis_author_seq \l_@@_thesis_tmp_tl
        \group_end:
      }
    \tl_gset:Ne \@author
      { \seq_use:Nn \g_@@_thesis_author_seq { \exp_not:N \and } }
  }
\RenewDocumentCommand \author { m } { \@@_thesis_author_parse:n {#1} }
%    \end{macrocode}
%
% Ein Feld eines Autors wird aus dessen Eigenschaftsliste gelesen,
% ist es dort nicht vorhanden, wird der globale Wert verwendet.
%    \begin{macrocode}
\cs_new_protected:Npn \@@_thesis_author_get:nnN #1#2#3
  {
    \prop_get:cnNF { g_@@_thesis_author_ #1 _prop } {#2} #3
      { \tl_set_eq:Nc #3 { @#2 } }
  }
%</body&peri>
%    \end{macrocode}
%
% \subsection{Lokalisierung}
%
% Die Bezeichner werden aus \dpkg{tudscr} übernommen und mittels \dmacro{\providecaptionname}
% für alle deutschen und englischen Sprachvarianten definiert.
% Sie lassen sich mit \dmacro{\renewcaptionname} überschreiben.
%
% Um zu prüfen, ob eine Liste von Personen mehr als eine Person enthält, wird ein Hilfsmakro bereitgestellt,
% welches auch in eigenen Bezeichnern genutzt werden kann.
%    \begin{macrocode}
%<*body&peri>
\cs_new_protected:Npn \tudcd@thesis@ifand #1
  { \tl_if_in:NnTF #1 { \and } }
%    \end{macrocode}
%
% Die Texte werden außerhalb der \dpkg{expl3}-Syntax definiert, um Leerzeichen setzen zu können.
%    \begin{macrocode}
\ExplSyntaxOff
\newcommand*\tudcd@thesis@german{%
  \providecaptionname{%
    german,ngerman,austrian,naustrian,swissgerman,nswissgerman%
  }%
}
\tudcd@thesis@german{\graduationtext}{%
  zur Erlangung des akademischen Grades%
}
\tudcd@thesis@german{\refereename}{%
  \tudcd@thesis@ifand{\@referee}{%
    \ifx\refereeothername\@empty Gutachter\else Erstgutachter\fi
  }{Gutachter}%
}
\tudcd@thesis@german{\refereeothername}{Zweitgutachter}
\tudcd@thesis@german{\advisorname}{%
  \tudcd@thesis@ifand{\@advisor}{%
    \ifx\advisorothername\@empty Fachreferenten\else Erster Fachreferent\fi
  }{Fachreferent}%
}
\tudcd@thesis@german{\advisorothername}{}
\tudcd@thesis@german{\supervisorname}{%
  \tudcd@thesis@ifand{\@supervisor}{%
    \ifx\supervisorothername\@empty Betreuer\else Erstbetreuer\fi
  }{Betreuer}%
}
\tudcd@thesis@german{\supervisorothername}{}
\tudcd@thesis@german{\professorname}{%
  \tudcd@thesis@ifand{\@professor}{%
    \ifx\professorothername\@empty
      Betreuende Hochschullehrer%
    \else
      Erster betreuender Hochschullehrer%
    \fi
  }{Betreuender Hochschullehrer}%
}
\tudcd@thesis@german{\professorothername}{}
\tudcd@thesis@german{\datetext}{Eingereicht am}
\tudcd@thesis@german{\dateofbirthtext}{Geboren am}
\tudcd@thesis@german{\placeofbirthtext}{in}
\tudcd@thesis@german{\defensedatetext}{Verteidigt am}
\tudcd@thesis@german{\matriculationyearname}{Immatrikulationsjahr}
\tudcd@thesis@german{\matriculationnumbername}{Matrikelnummer}
\tudcd@thesis@german{\coursename}{Studiengang}
\tudcd@thesis@german{\disciplinename}{Studienrichtung}
\tudcd@thesis@german{\abstractname}{Zusammenfassung}
\tudcd@thesis@german{\confirmationname}{Selbstst\"andigkeitserkl\"arung}
\tudcd@thesis@german{\confirmationtext}{%
  Hiermit versichere ich, dass ich das vorliegende Dokument %
  \ifx\tudcd@thesis@title\@empty\else
    mit dem Titel \emph{\tudcd@thesis@title} %
  \fi
  selbstst\"andig und ohne unzul\"assige Hilfe Dritter verfasst habe. %
  Es wurden keine anderen als die in diesem Dokument angegebenen Hilfsmittel %
  und Quellen benutzt. Die w\"ortlichen und sinngem\"a\ss{} %
  \"ubernommenen Zitate habe ich als solche kenntlich gemacht. %
  \ifx\@supporter\@empty
    Es waren keine weiteren Personen an der geistigen Herstellung %
    des vorliegenden Dokumentes beteiligt. %
  \else
    W\"ahrend der Anfertigung dieses Dokumentes wurde ich nur von %
    folgenden Personen unterst\"utzt:%
    \begin{quote}\def\and{\newline}\@supporter\end{quote}%
    \noindent Zus\"atzliche Personen waren an der geistigen Herstellung %
    des vorliegenden Dokumentes nicht beteiligt. %
  \fi
  Mir ist bekannt, dass die Nichteinhaltung dieser Erkl\"arung zum %
  nachtr\"aglichen Entzug des Hochschulabschlusses f\"uhren kann.%
}
\tudcd@thesis@german{\blockingname}{Sperrvermerk}
\tudcd@thesis@german{\blockingtext}{%
  Dieses Dokument %
  \ifx\tudcd@thesis@title\@empty\else
    mit dem Titel \emph{\tudcd@thesis@title} %
  \fi
  enth\"alt vertrauliche Informationen%
  \ifx\@company\@empty\else der Firma \emph{\@company}\fi. %
  Ver\"offentlichungen, Vervielf\"altigungen und Einsichtnahme~-- auch nur %
  auszugsweise~-- sind ohne ausdr\"uckliche Genehmigung %
  \ifx\@company\@empty\else der Firma \emph{\@company} \fi
  nicht gestattet. Es darf nur dem Betreuer an der Technischen Universit\"at %
  Dresden, den Gutachtern sowie den zust\"andigen Pr\"ufungsorganen und dem % 
  Pr\"ufungsamt sowie weiteren f\"ur das Pr\"ufungs-, Plagiatskontroll- oder %
  Rechtsschutzverfahren berechtigten Personen zug\"anglich gemacht werden.%
}
\tudcd@thesis@german{\habilitationname}{Habilitation}
\tudcd@thesis@german{\dissertationname}{Dissertation}
\tudcd@thesis@german{\diplomathesisname}{Diplomarbeit}
\tudcd@thesis@german{\masterthesisname}{Master-Arbeit}
\tudcd@thesis@german{\bachelorthesisname}{Bachelor-Arbeit}
\tudcd@thesis@german{\studentthesisname}{Studienarbeit}
\tudcd@thesis@german{\studentresearchname}{Gro\ss{}er Beleg}
\tudcd@thesis@german{\projectpapername}{Projektarbeit}
\tudcd@thesis@german{\seminarpapername}{Seminararbeit}
\tudcd@thesis@german{\termpapername}{Hausarbeit}
\tudcd@thesis@german{\researchname}{Forschungsbericht}
\tudcd@thesis@german{\logname}{Protokoll}
\tudcd@thesis@german{\internshipname}{Praktikumsbericht}
\tudcd@thesis@german{\reportname}{Bericht}
%    \end{macrocode}
%
%    \begin{macrocode}
\newcommand*\tudcd@thesis@english{%
  \providecaptionname{%
    american,australian,british,canadian,english,newzealand,UKenglish,USenglish%
  }%
}
\tudcd@thesis@english{\graduationtext}{to achieve the academic degree}
\tudcd@thesis@english{\refereename}{%
  \tudcd@thesis@ifand{\@referee}{%
    \ifx\refereeothername\@empty Referees\else First referee\fi
  }{Referee}%
}
\tudcd@thesis@english{\refereeothername}{Second referee}
\tudcd@thesis@english{\advisorname}{%
  \tudcd@thesis@ifand{\@advisor}{%
    \ifx\advisorothername\@empty Advisors\else First advisor\fi
  }{Advisor}%
}
\tudcd@thesis@english{\advisorothername}{}
\tudcd@thesis@english{\supervisorname}{%
  \tudcd@thesis@ifand{\@supervisor}{%
    \ifx\supervisorothername\@empty Supervisors\else First supervisor\fi
  }{Supervisor}%
}
\tudcd@thesis@english{\supervisorothername}{}
\tudcd@thesis@english{\professorname}{%
  \tudcd@thesis@ifand{\@professor}{%
    \ifx\professorothername\@empty
      Supervising professors%
    \else
      First supervising professor%
    \fi
  }{Supervising professor}%
}
\tudcd@thesis@english{\professorothername}{}
\tudcd@thesis@english{\datetext}{Submitted on}
\tudcd@thesis@english{\dateofbirthtext}{Born on}
\tudcd@thesis@english{\placeofbirthtext}{in}
\tudcd@thesis@english{\defensedatetext}{Defended on}
\tudcd@thesis@english{\matriculationyearname}{Matriculation year}
\tudcd@thesis@english{\matriculationnumbername}{Matriculation number}
\tudcd@thesis@english{\coursename}{Course}
\tudcd@thesis@english{\disciplinename}{Discipline}
\tudcd@thesis@english{\abstractname}{Abstract}
\tudcd@thesis@english{\confirmationname}{Statement of authorship}
\tudcd@thesis@english{\confirmationtext}{%
  I hereby certify that I have authored this document %
  \ifx\tudcd@thesis@title\@empty\else
    entitled \emph{\tudcd@thesis@title} %
  \fi
  independently and without undue assistance from third %
  parties. No other than the resources and references %
  indicated in this document have been used. I have marked %
  both literal and accordingly adopted quotations as such. %
  \ifx\@supporter\@empty
    There were no additional persons involved in the %
  \else
    During the preparation of this document I was only %
    supported by the following persons:%
    \begin{quote}\def\and{\newline}\@supporter\end{quote}%
    \noindent Additional persons were not involved in the %
  \fi
  intellectual preparation of the present document. %
  I am aware that violations of this declaration may lead to %
  subsequent withdrawal of the academic degree.%
}
\tudcd@thesis@english{\blockingname}{Restriction note}
\tudcd@thesis@english{\blockingtext}{%
  This document %
  \ifx\tudcd@thesis@title\@empty\else
    entitled \emph{\tudcd@thesis@title} %
  \fi
  contains confidential data%
  \ifx\@company\@empty\else, disclosed by \emph{\@company}\fi. %
  Publications, duplications and inspections---even in part---are prohibited %
  without explicit permission%
  \ifx\@company\@empty\else\space by \emph{\@company}\fi, %
  as well as publications about the content of this document. %
  It may only be made accessible to the supervisor at Technische Universit\"at %
  Dresden, the reviewers and also the members of the examination board.%
}
\tudcd@thesis@english{\habilitationname}{Habilitation}
\tudcd@thesis@english{\dissertationname}{Dissertation}
\tudcd@thesis@english{\diplomathesisname}{Diploma Thesis}
\tudcd@thesis@english{\masterthesisname}{Master Thesis}
\tudcd@thesis@english{\bachelorthesisname}{Bachelor Thesis}
\tudcd@thesis@english{\studentthesisname}{Student Thesis}
\tudcd@thesis@english{\studentresearchname}{Student Research Project}
\tudcd@thesis@english{\projectpapername}{Project Paper}
\tudcd@thesis@english{\seminarpapername}{Seminar Paper}
\tudcd@thesis@english{\termpapername}{Term Paper}
\tudcd@thesis@english{\researchname}{Research Report}
\tudcd@thesis@english{\logname}{Log}
\tudcd@thesis@english{\internshipname}{Internship Report}
\tudcd@thesis@english{\reportname}{Report}
\ExplSyntaxOn
%</body&peri>
%    \end{macrocode}
%
% \subsection{Die Titelseite}
%
% \begin{codedescribe}[macro,new=0.8.0]{\maketitle}
% \begin{codesyntax}
% \tsmacro{\maketitle}[Seitenzahl]
% \end{codesyntax}
% Setzt die Titelseite mit Logo, Organisationseinheiten, Art der Arbeit, Titel, Autoren, akademischem Grad,
% den betreuenden Personen sowie den Daten der Einreichung und Verteidigung.
% Anschließend folgen, sofern angegeben, die Titelrückseite (\dmacro{\uppertitleback}, \dmacro{\lowertitleback})
% und die Widmung (\dmacro{\dedication}).
% \end{codedescribe}
%
% Eine Zeile mit Angaben wird nur ausgegeben, wenn der Inhalt nicht leer ist.
%    \begin{macrocode}
%<*body&peri>
\dim_new:N \l_@@_thesis_title_logo_dim
\dim_set:Nn \l_@@_thesis_title_logo_dim { 15mm }
\cs_new_protected:Npn \@@_thesis_title_line:N #1
  { \tl_if_empty:NF #1 { #1 \par } }
\cs_new_protected:Npn \@@_thesis_title_line:nN #1#2
  { \tl_if_empty:NF #2 { #1 : ~ #2 \par } }
%    \end{macrocode}
%
% Der Block eines Autors besteht aus dem Namen und den zugehörigen Angaben.
%    \begin{macrocode}
\tl_new:N \l_@@_thesis_title_a_tl
\tl_new:N \l_@@_thesis_title_b_tl
\cs_new_protected:Npn \@@_thesis_title_author:nn #1#2
  {
    { \usekomafont{author} #2 \par }
    \smallskip 
    \group_begin:
      \usekomafont{titleinfo}
      \@@_thesis_author_get:nnN {#1} { emailaddress } \l_@@_thesis_title_a_tl
      \tl_if_empty:NF \l_@@_thesis_title_a_tl
        {
          \prop_if_in:cnTF { g_@@_thesis_author_ #1 _prop } { emailaddressnolink }
            { \bool_set_true:N \l_tmpa_bool }
            { \bool_set_eq:NN \l_tmpa_bool \g_@@_thesis_emailaddress_nolink_bool }
          \bool_lazy_or:nnTF { \l_tmpa_bool } { ! \cs_if_exist_p:N \href }
            { \l_@@_thesis_title_a_tl \par }
            {
              \exp_args:Ne \href { mailto: \l_@@_thesis_title_a_tl }
                { \l_@@_thesis_title_a_tl } \par
            }
        }
      \@@_thesis_author_get:nnN {#1} { authormore } \l_@@_thesis_title_a_tl
      \@@_thesis_title_line:N \l_@@_thesis_title_a_tl
      \@@_thesis_author_get:nnN {#1} { course } \l_@@_thesis_title_a_tl
      \@@_thesis_title_line:nN { \coursename } \l_@@_thesis_title_a_tl
      \@@_thesis_author_get:nnN {#1} { discipline } \l_@@_thesis_title_a_tl
      \@@_thesis_title_line:nN { \disciplinename } \l_@@_thesis_title_a_tl
      \@@_thesis_author_get:nnN {#1} { dateofbirth } \l_@@_thesis_title_a_tl
      \@@_thesis_author_get:nnN {#1} { placeofbirth } \l_@@_thesis_title_b_tl
      \tl_if_empty:NF \l_@@_thesis_title_a_tl
        {
          \dateofbirthtext \c_space_tl \l_@@_thesis_title_a_tl
          \tl_if_empty:NF \l_@@_thesis_title_b_tl
            { \c_space_tl \placeofbirthtext \c_space_tl \l_@@_thesis_title_b_tl }
          \par
        }
      \@@_thesis_author_get:nnN {#1} { matriculationnumber } \l_@@_thesis_title_a_tl
      \@@_thesis_title_line:nN { \matriculationnumbername } \l_@@_thesis_title_a_tl
      \@@_thesis_author_get:nnN {#1} { matriculationyear } \l_@@_thesis_title_a_tl
      \@@_thesis_title_line:nN { \matriculationyearname } \l_@@_thesis_title_a_tl
    \group_end:
  }
\cs_new_protected:Npn \@@_thesis_title_authors:
  {
    \seq_if_empty:NTF \g_@@_thesis_author_seq
      { { \usekomafont{author} \@author \par } }
      {
        \seq_map_indexed_inline:Nn \g_@@_thesis_author_seq
          {
            \int_compare:nNnF {##1} = { 1 } { \medskip }
            \@@_thesis_title_author:nn {##1} {##2}
          }
      }
  }
%    \end{macrocode}
%
% Die betreuenden Personen werden in einer Tabelle ausgegeben.
% Enthält eine Liste mehrere Personen und ist die Bezeichnung für die weiteren Personen (bspw.\@ \dmacro{\refereeothername})
% nicht leer, so erhält jede weitere Person eine eigene Bezeichnung, ansonsten wird die Mehrzahl verwendet.
%    \begin{macrocode}
\tl_new:N \l_@@_thesis_title_rows_tl
\cs_generate_variant:Nn \seq_set_split:Nnn { Nnv }
\cs_new_protected:Npn \@@_thesis_title_row:nn #1#2
  {
    \tl_put_right:Nn \l_@@_thesis_title_rows_tl
      { \tl_if_blank:nF {#1} { #1 : } & #2 \\ }
  }
\cs_new_protected:Npn \@@_thesis_title_names:n #1
  {
    \tl_if_empty:cF { @#1 }
      {
        \seq_set_split:Nnv \l_@@_thesis_tmp_seq { \and } { @#1 }
        \bool_lazy_and:nnTF
          { \int_compare_p:nNn { \seq_count:N \l_@@_thesis_tmp_seq } > { 1 } }
          { ! \tl_if_empty_p:c { #1othername } }
          {
            \seq_map_indexed_inline:Nn \l_@@_thesis_tmp_seq
              {
                \int_compare:nNnTF {##1} = { 1 }
                  { \@@_thesis_title_row:nn { \use:c { #1name } } {##2} }
                  { \@@_thesis_title_row:nn { \use:c { #1othername } } {##2} }
              }
          }
          {
            \seq_map_indexed_inline:Nn \l_@@_thesis_tmp_seq
              {
                \int_compare:nNnTF {##1} = { 1 }
                  { \@@_thesis_title_row:nn { \use:c { #1name } } {##2} }
                  { \@@_thesis_title_row:nn { } {##2} }
              }
          }
      }
  }
%    \end{macrocode}
%
% Der eigentliche Satz der Titelseite.
%    \begin{macrocode}
\cs_new_protected:Npn \@@_thesis_title_page:
  {
    \group_begin:
    \normalfont \normalsize
    \setparsizes { 0pt } { 0pt } { 0pt plus 1fil }
    \raggedright
    \noindent \TypesetTUDLogo { \l_@@_thesis_title_logo_dim } \par
    \vspace { 2 \baselineskip }
    \group_begin:
      \usekomafont{titleinfo}
      \@@_thesis_title_line:N \@faculty
      \@@_thesis_title_line:N \@department
      \@@_thesis_title_line:N \@institute
      \@@_thesis_title_line:N \@chair
    \group_end:
    \vspace { 0pt plus 1fill }
    \tl_if_empty:NF \@titlehead { { \usekomafont{titlehead} \@titlehead \par } \bigskip }
    \bool_lazy_or:nnF
      { \tl_if_empty_p:N \@subject }
      { \tl_if_eq_p:NN \@subject \@thesis }
      { { \usekomafont{subject} \@subject \par } \medskip }
    \tl_if_empty:NF \@thesis { { \usekomafont{thesis} \@thesis \par } \bigskip }
    { \usekomafont{title} \@title \par }
    \tl_if_empty:NF \@subtitle { \medskip { \usekomafont{subtitle} \@subtitle \par } }
    \vspace { 3 \baselineskip }
    \@@_thesis_title_authors:
    \vspace { 0pt plus 1fill }
    \tl_if_empty:NF \@graduation
      {
        { \usekomafont{titleinfo} \graduationtext \par }
        \smallskip
        {
          \usekomafont{thesis} \@graduation
          \tl_if_empty:NF \@graduationabbreviation { \nobreakspace ( \@graduationabbreviation ) }
          \par
        }
        \vspace { 0pt plus 1fill }
      }
    \tl_clear:N \l_@@_thesis_title_rows_tl
    \@@_thesis_title_names:n { referee }
    \@@_thesis_title_names:n { advisor }
    \@@_thesis_title_names:n { supervisor }
    \@@_thesis_title_names:n { professor }
    \tl_if_empty:NF \@thesis
      {
        \tl_if_empty:NF \@date
          { \@@_thesis_title_row:nn { \datetext } { \usekomafont{date} \@date } }
        \tl_if_empty:NF \@defensedate
          { \@@_thesis_title_row:nn { \defensedatetext } { \usekomafont{date} \@defensedate } }
      }
    \tl_if_empty:NF \l_@@_thesis_title_rows_tl
      {
        { \usekomafont{titleinfo}
          \begin{tabular} { @{} l @{\hspace{1em}} l @{} }
            \l_@@_thesis_title_rows_tl
          \end{tabular} \par }
      }
    \tl_if_empty:NT \@thesis
      { \tl_if_empty:NF \@date { { \usekomafont{date} \@date \par } } }
    \tl_if_empty:NF \@publishers
      { \bigskip { \usekomafont{publishers} \@publishers \par } }
    \@thanks
    \group_end:
  }
%    \end{macrocode}
%
% \dmacro{\maketitle} setzt den Schmutztitel, die Titelseite, die Titelrückseite und die Widmung jeweils auf eigene Seiten.
%    \begin{macrocode}
\RenewDocumentCommand \maketitle { o }
  {
    \group_begin:
    \cs_set_eq:NN \footnote \thanks
    \tl_if_empty:NF \@extratitle
      {
        \cleardoubleoddemptypage
        \thispagestyle{empty}
        \group_begin: \noindent \@extratitle \par \group_end:
        \cleardoubleoddemptypage
      }
    \clearpage
    \IfValueT {#1} { \setcounter { page } {#1} }
    \thispagestyle{empty}
    \@@_thesis_title_page:
    \bool_lazy_and:nnT
      { \tl_if_empty_p:N \@uppertitleback }
      { \tl_if_empty_p:N \@lowertitleback }
      { \legacy_if:nF { @twoside } { \clearpage } }
    \bool_lazy_and:nnF
      { \tl_if_empty_p:N \@uppertitleback }
      { \tl_if_empty_p:N \@lowertitleback }
      {
        \clearpage
        \thispagestyle{empty}
        \group_begin:
          \normalfont \normalsize
          \noindent \begin{minipage}[t]{\textwidth} \@uppertitleback \end{minipage} \par
          \vfill
          \noindent \begin{minipage}[b]{\textwidth} \@lowertitleback \end{minipage} \par
        \group_end:
        \clearpage
      }
    \tl_if_empty:NF \@dedication
      {
        \cleardoubleoddemptypage
        \thispagestyle{empty}
        \null \vfill
        { \centering \usekomafont{dedication} \@dedication \par }
        \vspace { 0pt plus 3fill }
        \clearpage
      }
    \cleardoubleoddemptypage
    \group_end:
    \cs_gset_eq:NN \thanks \use_none:n
  }
%</body&peri>
%    \end{macrocode}
%
% \subsection{Zusammenfassung und Erklärungen}
%
% Beide Umgebungen werden durch eine gemeinsame Maschinerie gesetzt, welche die Logik von \dpkg{tudscr} nachbildet.
% Ein \emph{Teil} ist eine einzelne Zusammenfassung (getrennt durch \dmacro{\nextabstract}) oder
% eine einzelne Erklärung.
%
% \subsubsection{Parameter}
%
% Die optionalen Argumente der Umgebungen und Befehle sind Schlüssel-Wert-Listen.
% Enthält das Argument kein Gleichheitszeichen, so wird es als Wert eines Standardschlüssels interpretiert
% (Sprache bzw.\@ bei einer Zahl die Spaltenanzahl, bei \dmacro{\confirmation} unterstützende Personen
% und bei \dmacro{\blocking} das Unternehmen).
%    \begin{macrocode}
%<*body&peri>
\tl_new:N \l_@@_thesis_fm_kind_tl
\tl_new:N \l_@@_thesis_fm_ps_tl
\tl_new:N \l_@@_thesis_fm_oldps_tl
\tl_new:N \l_@@_thesis_fm_partopts_tl
\tl_new:N \l_@@_thesis_fm_bare_tl
\int_new:N \l_@@_thesis_fm_columns_int
\bool_new:N \l_@@_thesis_fm_titlepage_bool
\bool_new:N \g_@@_thesis_fm_first_bool
\bool_new:N \l_@@_thesis_in_declarations_bool
\msg_new:nnn { tudcd } { thesis/fm-pagestyle }
  { The~page~style~'#1'~is~not~defined. }
\msg_new:nnn { tudcd } { thesis/fm-pagestyle-ignored }
  {
    The~parameter~'pagestyle'~is~ignored,~because~'#1'~is~not~set~on~separate~pages.~
    Use~'#1=chapter'~or~the~KOMA-Script~option~'titlepage'.
  }
\msg_new:nnn { tudcd } { thesis/fm-multicol }
  { The~parameter~'columns'~requires~the~package~'multicol'. }
\msg_new:nnn { tudcd } { thesis/fm-language }
  { The~parameter~'language'~requires~a~language~package~like~'babel'. }
\msg_new:nnn { tudcd } { thesis/fm-tocmultiple }
  { 'tocmultiple'~cannot~be~used~with~'#1=single'~or~'#1=chapter'. }
\msg_new:nnn { tudcd } { thesis/fm-usage }
  { '\iow_char:N\\#1'~can~only~be~used~within~the~environment~'#2'. }
\keys_define:nn { tudcd / thesis / fm }
  {
    language .code:n =
      {
        \cs_if_exist:NTF \selectlanguage
          { \selectlanguage {#1} }
          { \msg_error:nn { tudcd } { thesis/fm-language } }
      },
    columns .int_set:N = \l_@@_thesis_fm_columns_int,
    pagestyle .code:n =
      {
        \bool_if:NTF \l_@@_thesis_fm_titlepage_bool
          {
            \cs_if_exist:cTF { ps@#1 }
              { \tl_set:Nn \l_@@_thesis_fm_ps_tl {#1} }
              { \msg_error:nnn { tudcd } { thesis/fm-pagestyle } {#1} }
          }
          { \msg_warning:nnV { tudcd } { thesis/fm-pagestyle-ignored } \l_@@_thesis_fm_kind_tl }
      },
    markboth .code:n =
      {
        \str_case:nnF {#1}
          {
            { true }
              {
                \bool_set_true:c { l_@@_thesis_ \l_@@_thesis_fm_kind_tl _markboth_bool }
                \tl_if_empty:NT \l_@@_thesis_fm_ps_tl { \tl_set:Nn \l_@@_thesis_fm_ps_tl { headings } }
              }
            { false }
              { \bool_set_false:c { l_@@_thesis_ \l_@@_thesis_fm_kind_tl _markboth_bool } }
          }
          {
            \bool_set_false:c { l_@@_thesis_ \l_@@_thesis_fm_kind_tl _markboth_bool }
            \tl_if_empty:NT \l_@@_thesis_fm_ps_tl { \tl_set:Nn \l_@@_thesis_fm_ps_tl { headings } }
            \manualmark
            \markboth {#1} {#1}
          }
      },
    markboth .default:n = true,
    option .code:n = { \@@_thesis_fm_set:Vn \l_@@_thesis_fm_kind_tl {#1} },
    abstract .meta:n = { option = {#1} },
    declaration .meta:n = { option = {#1} },
    company .code:n = { \tl_set:Nn \@company {#1} },
    supporter .code:n = { \tl_set:Nn \@supporter {#1} },
    place .code:n = { \tl_set:Nn \@place {#1} },
    closing .code:n = { \tl_set:Nn \@confirmationclosing {#1} },
    confirmationclosing .meta:n = { closing = {#1} },
    unknown .code:n = { \@@_thesis_fm_bare:e { \l_keys_key_str } },
  }
\cs_generate_variant:Nn \@@_thesis_fm_set:nn { V }
\cs_new_protected:Npn \@@_thesis_fm_bare:n #1
  {
    \tl_if_empty:NTF \l_@@_thesis_fm_bare_tl
      {
        \regex_match:nnTF { \A \d+ \Z } {#1}
          { \int_set:Nn \l_@@_thesis_fm_columns_int {#1} }
          { \keys_set:nn { tudcd / thesis / fm } { language = {#1} } }
      }
      {
        \use:e
          {
            \keys_set:nn { tudcd / thesis / fm }
              { \l_@@_thesis_fm_bare_tl = { \exp_not:n {#1} } }
          }
      }
  }
\cs_generate_variant:Nn \@@_thesis_fm_bare:n { e }
\cs_new_protected:Npn \@@_thesis_fm_params:n #1
  {
    \tl_if_in:nnTF {#1} { = }
      { \keys_set:nn { tudcd / thesis / fm } {#1} }
      { \tl_if_blank:nF {#1} { \@@_thesis_fm_bare:n {#1} } }
  }
%    \end{macrocode}
%
% \subsubsection{Überschriften}
%
% Die Überschrift eines Teils hängt von der Ebene ab.
% Anschließend werden gegebenenfalls Kolumnentitel und Einträge ins Inhaltsverzeichnis gesetzt.
%    \begin{macrocode}
\cs_new_protected:Npn \@@_thesis_fm_head:n #1
  {
    \int_case:nn { \int_use:c { l_@@_thesis_ \l_@@_thesis_fm_kind_tl _level_int } }
      {
        { 0 }
          {
            \bool_if:NT \l_@@_thesis_fm_titlepage_bool { \par }
            \cs_if_exist_use:N \phantomsection
            \vskip \bigskipamount
            \vskip -\parskip
          }
        { 1 }
          {
            \bool_if:NT \l_@@_thesis_fm_titlepage_bool { \par }
            \cs_if_exist_use:N \phantomsection
            \vskip \bigskipamount
            \vskip -\parskip
            \group_begin:
              \noindent \centering
              \normalfont \sectfont \nobreak #1
              \par
            \group_end:
            \nopagebreak
            \vskip \medskipamount
            \vskip -\parskip
            \@afterheading
          }
        { 2 }
          {
            \section* {#1}
            \cs_if_eq:NNF \@mkboth \use_none:nn { \markright { \MakeMarkcase {#1} } }
          }
        { 3 }
          {
            \chapter* {#1}
            \markboth { \MakeMarkcase {#1} } { \MakeMarkcase {#1} }
          }
      }
    \bool_if:cT { l_@@_thesis_ \l_@@_thesis_fm_kind_tl _markboth_bool }
      { \markboth { \MakeMarkcase {#1} } { \MakeMarkcase {#1} } }
    \int_case:nn { \int_use:c { l_@@_thesis_ \l_@@_thesis_fm_kind_tl _toc_int } }
      {
        { 1 } { \cs_if_exist_use:N \phantomsection \addchaptertocentry { } {#1} }
        { 2 } { \cs_if_exist_use:N \phantomsection \addsectiontocentry { } {#1} }
      }
  }
%    \end{macrocode}
%
% \subsubsection{Teile}
%
% Vor jedem Teil wird der vertikale Abstand bzw.\@ der Seitenumbruch gesetzt.
% Werden mehrere Teile gemeinsam auf einer Seite gesetzt (\dcode{multiple}), so werden sie
% bei \dcode{fil} gleichmäßig vertikal verteilt, ansonsten beginnt jeder Teil auf einer neuen Seite.
%    \begin{macrocode}
\cs_new_protected:Npn \@@_thesis_fm_separate:
  {
    \legacy_if:nF { @twocolumn }
      {
        \bool_if:NT \l_@@_thesis_fm_titlepage_bool
          {
            \bool_if:NTF \g_@@_thesis_fm_first_bool
              {
                \bool_if:cT { l_@@_thesis_ \l_@@_thesis_fm_kind_tl _fil_bool }
                  { \null \vfil }
              }
              {
                \bool_if:cTF { l_@@_thesis_ \l_@@_thesis_fm_kind_tl _multi_bool }
                  {
                    \bool_if:cT { l_@@_thesis_ \l_@@_thesis_fm_kind_tl _fil_bool }
                      { \vfil }
                  }
                  {
                    \bool_if:cTF { l_@@_thesis_ \l_@@_thesis_fm_kind_tl _fil_bool }
                      { \vfil \null \clearpage \null \vfil }
                      { \clearpage }
                  }
              }
          }
      }
    \bool_gset_false:N \g_@@_thesis_fm_first_bool
  }
%    \end{macrocode}
% Ein Teil wird in Zitatform (\denv{quotation}, bzw.\@ \denv{quoting} falls geladen) gesetzt, wenn er
% weder auf eigener Seite steht noch eine Gliederungsüberschrift besitzt. Bei mehreren Spalten wird die
% Überschrift über die Spalten gesetzt.
%    \begin{macrocode}
\tl_new:N \l_@@_thesis_fm_quoting_tl
\tl_set:Nn \l_@@_thesis_fm_quoting_tl { quotation }
\hook_gput_code:nnn { package/quoting/after } { tudcd/thesis }
  { \tl_set:Nn \l_@@_thesis_fm_quoting_tl { quoting } }
\cs_new_protected:Npn \@@_thesis_fm_part_start:nn #1#2
  {
    \group_begin:
      \@@_thesis_fm_params:n {#1}
      \@@_thesis_fm_separate:
      \bool_lazy_and:nnT
        { ! \l_@@_thesis_fm_titlepage_bool }
        { \int_compare_p:nNn { \int_use:c { l_@@_thesis_ \l_@@_thesis_fm_kind_tl _level_int } } < { 2 } }
        { \legacy_if:nF { @twocolumn } { \exp_args:NV \begin \l_@@_thesis_fm_quoting_tl } }
      \int_compare:nNnTF \l_@@_thesis_fm_columns_int > { 1 }
        {
          \vspace { -\multicolsep }
          \begin{multicols} { \l_@@_thesis_fm_columns_int } [ \@@_thesis_fm_head:n {#2} ]
        }
        { \@@_thesis_fm_head:n {#2} }
  }
\cs_new_protected:Npn \@@_thesis_fm_part_end:
  {
      \int_compare:nNnTF \l_@@_thesis_fm_columns_int > { 1 }
        {
          \end{multicols}
          \vspace { -\multicolsep }
        }
        { \par }
      \bool_lazy_and:nnT
        { ! \l_@@_thesis_fm_titlepage_bool }
        { \int_compare_p:nNn { \int_use:c { l_@@_thesis_ \l_@@_thesis_fm_kind_tl _level_int } } < { 2 } }
        { \legacy_if:nF { @twocolumn } { \exp_args:NV \end \l_@@_thesis_fm_quoting_tl \par } }
    \group_end:
  }
\cs_new_protected:Npn \@@_thesis_fm_part:nnn #1#2#3
  {
    \@@_thesis_fm_part_start:nn {#1} {#2}
      #3
    \@@_thesis_fm_part_end:
  }
%    \end{macrocode}
%
% \subsubsection{Die gemeinsame Umgebung}
%
% Die Umgebungen setzen den Inhalt \dcode{\#3} der Art \dcode{\#1} mit den Parametern \dcode{\#2}.
% Die Ebene \dcode{chapter} erzwingt eigene Seiten, auf welchen jeder Teil einzeln und ohne vertikale Zentrierung steht.
% Bei \dcode{tocmultiple} wird der Inhalt vorab in einer Box gesetzt, um die Überschriften aller Teile
% für einen gemeinsamen Eintrag im Inhaltsverzeichnis zu sammeln.
%    \begin{macrocode}
\tl_new:N \g_@@_thesis_fm_toc_tl
\cs_new_protected:Npn \@@_thesis_fm_env:nnn #1#2#3
  {
    \group_begin:
    \tl_set:Nn \l_@@_thesis_fm_kind_tl {#1}
    \tl_clear:N \l_@@_thesis_fm_ps_tl
    \int_set:Nn \l_@@_thesis_fm_columns_int { 1 }
    \legacy_if:nTF { @titlepage }
      { \bool_set_true:N \l_@@_thesis_fm_titlepage_bool }
      { \bool_set_false:N \l_@@_thesis_fm_titlepage_bool }
    \int_compare:nNnT { \int_use:c { l_@@_thesis_#1_level_int } } = { 3 }
      {
        \bool_set_true:N \l_@@_thesis_fm_titlepage_bool
        \bool_set_false:c { l_@@_thesis_#1_multi_bool }
        \bool_set_false:c { l_@@_thesis_#1_fil_bool }
      }
    \@@_thesis_fm_params:n {#2}
    \bool_lazy_and:nnT
      { \int_compare_p:nNn \l_@@_thesis_fm_columns_int > { 1 } }
      { ! \cs_if_exist_p:N \multicols }
      {
        \msg_error:nn { tudcd } { thesis/fm-multicol }
        \int_set:Nn \l_@@_thesis_fm_columns_int { 1 }
      }
    \bool_if:NT \l_@@_thesis_fm_titlepage_bool
      {
        \clearpage
        \tl_set:Ne \l_@@_thesis_fm_oldps_tl { \currentpagestyle }
        \tl_if_empty:NTF \l_@@_thesis_fm_ps_tl
          { \pagestyle { empty } }
          { \pagestyle { \l_@@_thesis_fm_ps_tl } }
      }
    \int_compare:nNnT { \int_use:c { l_@@_thesis_#1_toc_int } } > { 2 }
      { \@@_thesis_fm_toc_multiple:nn {#1} {#3} }
    \bool_gset_true:N \g_@@_thesis_fm_first_bool
    #3
    \bool_if:NTF \l_@@_thesis_fm_titlepage_bool
      {
        \bool_if:cT { l_@@_thesis_#1_fil_bool } { \vfil \null }
        \clearpage
        \exp_args:NV \pagestyle \l_@@_thesis_fm_oldps_tl
      }
      {
        \int_compare:nNnT { \int_use:c { l_@@_thesis_#1_level_int } } = { 2 }
          { \vskip \bigskipamount \vskip -\parskip }
      }
    \group_end:
  }
\cs_new_protected:Npn \@@_thesis_fm_toc_multiple:nn #1#2
  {
    \bool_if:cTF { l_@@_thesis_#1_multi_bool }
      {
        \tl_gclear:N \g_@@_thesis_fm_toc_tl
        \group_begin:
          \cs_set_eq:NN \clearpage \scan_stop:
          \cs_set_eq:NN \addtocentrydefault \use_none:nnn
          \cs_set_protected:Npn \@@_thesis_fm_head:n ##1
            {
              \protected@xdef \g_@@_thesis_fm_toc_tl
                {
                  \tl_if_empty:NF \g_@@_thesis_fm_toc_tl { \g_@@_thesis_fm_toc_tl / }
                  ##1
                }
            }
          \vbox_set:Nn \l_@@_thesis_tmp_box {#2}
        \group_end:
        \tl_if_empty:NF \g_@@_thesis_fm_toc_tl
          {
            \cs_if_exist_use:N \phantomsection
            \int_compare:nNnTF { \int_use:c { l_@@_thesis_#1_toc_int } } = { 3 }
              { \exp_args:NnV \addchaptertocentry { } \g_@@_thesis_fm_toc_tl }
              { \exp_args:NnV \addsectiontocentry { } \g_@@_thesis_fm_toc_tl }
          }
      }
      {
        \msg_warning:nnn { tudcd } { thesis/fm-tocmultiple } {#1}
        \int_sub:cn { l_@@_thesis_#1_toc_int } { 2 }
      }
  }
%    \end{macrocode}
%
% \subsubsection{Die Zusammenfassung}
%
% \begin{codedescribe}[env,new=0.8.0]{abstract}
% \begin{codesyntax}
% \tsmacro{\begin{abstract}}[Parameter]
% \ldots
% \tsmacro{\nextabstract}[Parameter]
% \ldots
% \tsmacro{\end{abstract}}{}
% \end{codesyntax}
% Setzt eine oder mehrere Zusammenfassungen. Mit \dmacro{\nextabstract} wird die nächste Zusammenfassung,
% üblicherweise in einer anderen Sprache, begonnen. Mögliche Parameter sind
% \dcode{language}, \dcode{columns}, \dcode{pagestyle}, \dcode{markboth} und \dcode{option}
% (bzw.\@ \dcode{abstract}), welches die Werte der Klassenoption \doption{abstract} lokal setzt.
% \end{codedescribe}
%    \begin{macrocode}
\NewDocumentCommand \nextabstract { O{} }
  { \msg_error:nnnn { tudcd } { thesis/fm-usage } { nextabstract } { abstract } }
\DeclareDocumentEnvironment { abstract } { O{} +b }
  {
    \@@_thesis_fm_env:nnn { abstract } {#1}
      {
        \RenewDocumentCommand \nextabstract { O{} }
          {
            \@@_thesis_fm_part_end:
            \@@_thesis_fm_part_start:nn {##1} { \abstractname }
          }
        \@@_thesis_fm_part:nnn {#1} { \abstractname } {#2}
      }
  }
  { }
%    \end{macrocode}
%
% \subsubsection{Die Erklärungen}
%
% \begin{codedescribe}[env,new=0.8.0]{declarations}
% \begin{codesyntax}
% \tsmacro{\begin{declarations}}[Parameter]
% \ldots
% \tsmacro{\end{declarations}}{}
% \end{codesyntax}
% Fasst mehrere Erklärungen zusammen. Neben den Parametern der Zusammenfassung sind
% \dcode{company}, \dcode{supporter}, \dcode{place} und \dcode{closing} möglich.
% \end{codedescribe}
%    \begin{macrocode}
\DeclareDocumentEnvironment { declarations } { O{} +b }
  {
    \@@_thesis_fm_env:nnn { declaration } {#1}
      {
        \bool_set_true:N \l_@@_thesis_in_declarations_bool
        \cs_set:Npn \and { , ~ }
        \cs_set_eq:NN \thanks \use_none:n
        #2
      }
  }
  { }
%    \end{macrocode}
%
% Wird ein Befehl für eine Erklärung außerhalb der Umgebung genutzt, so wird diese implizit geöffnet.
%    \begin{macrocode}
\cs_new_protected:Npn \@@_thesis_declaration:nn #1#2
  {
    \bool_if:NTF \l_@@_thesis_in_declarations_bool
      {
        \tl_set:Nn \l_@@_thesis_fm_partopts_tl {#1}
        #2
      }
      {
        \tl_clear:N \l_@@_thesis_fm_partopts_tl
        \begin{declarations} [#1] #2 \end{declarations}
      }
  }
\cs_new_protected:Npn \@@_thesis_declaration_part:nn #1#2
  { \exp_args:NV \@@_thesis_fm_part:nnn \l_@@_thesis_fm_partopts_tl {#1} {#2} }
%    \end{macrocode}
%
% \begin{codedescribe}[macro,new=0.8.0]{\confirmation,\blocking,\declaration,\nextdeclaration}
% \begin{codesyntax}
% \tsmacro{\confirmation}[Parameter]
%
% \tsmacro{\blocking}[Parameter]
%
% \tsmacro{\declaration}[Parameter]
%
% \tsmacro{\nextdeclaration}[Parameter]{Überschrift,Text}
% \end{codesyntax}
% \dmacro{\confirmation} setzt die Selbstständigkeitserklärung, \dmacro{\blocking} den Sperrvermerk und
% \dmacro{\declaration} beide. Bei den gesternten Varianten werden die Betreuenden als unterstützende Personen genannt.
% \dmacro{\nextdeclaration} setzt eine eigene Erklärung und ist nur innerhalb von \denv{declarations} erlaubt.
% \end{codedescribe}
%
% Der Abschluss der Selbstständigkeitserklärung besteht aus Ort, Datum und je einer Unterschriftenzeile pro Autor.
%    \begin{macrocode}
\tl_gset:Nn \@confirmationclosing { \@@_thesis_confirmation_closing: }
\cs_new_protected:Npn \@@_thesis_confirmation_closing:
  {
    \tl_if_empty:NF \@date
      {
        \par \medskip \noindent
        \tl_if_empty:NF \@place { \@place , \nobreakspace }
        \@date
      }
    \par
    \vspace { 3 \baselineskip }
    \noindent
    \seq_if_empty:NTF \g_@@_thesis_author_seq
      { \@@_thesis_signature:n { \@author } }
      {
        \seq_map_indexed_inline:Nn \g_@@_thesis_author_seq
          {
            \int_compare:nNnF {##1} = { 1 } { \hfill }
            \@@_thesis_signature:n {##2}
          }
      }
    \par
  }
\cs_new_protected:Npn \@@_thesis_signature:n #1
  {
    \parbox[t] { 0.4 \linewidth }
      { \rule { \linewidth } { 0.4pt } \\ \small #1 }
  }
\cs_new_protected:Npn \@@_thesis_declaration_opts:nn #1#2
  {
    \tl_set:Nn \l_@@_thesis_tmp_tl {#2}
    \bool_lazy_and:nnT {#1} { ! \tl_if_empty_p:N \@supervisor }
      { \tl_put_left:Ne \l_@@_thesis_tmp_tl { supporter = { \exp_not:o \@supervisor } , } }
  }
\cs_new_protected:Npn \@@_thesis_confirmation_part:
  {
    \@@_thesis_declaration_part:nn { \confirmationname }
      { \confirmationtext \par \@confirmationclosing }
  }
\cs_new_protected:Npn \@@_thesis_blocking_part:
  { \@@_thesis_declaration_part:nn { \blockingname } { \blockingtext } }
\NewDocumentCommand \confirmation { s O{} }
  {
    \group_begin:
      \tl_set:Nn \l_@@_thesis_fm_bare_tl { supporter }
      \@@_thesis_declaration_opts:nn {#1} {#2}
      \exp_args:NV \@@_thesis_declaration:nn \l_@@_thesis_tmp_tl
        { \@@_thesis_confirmation_part: }
    \group_end:
  }
\NewDocumentCommand \blocking { O{} }
  {
    \group_begin:
      \tl_set:Nn \l_@@_thesis_fm_bare_tl { company }
      \@@_thesis_declaration:nn {#1} { \@@_thesis_blocking_part: }
    \group_end:
  }
\NewDocumentCommand \declaration { s O{} }
  {
    \group_begin:
      \@@_thesis_declaration_opts:nn {#1} {#2}
      \exp_args:NV \@@_thesis_declaration:nn \l_@@_thesis_tmp_tl
        {
          \@@_thesis_confirmation_part:
          \@@_thesis_blocking_part:
        }
    \group_end:
  }
\NewDocumentCommand \nextdeclaration { O{} m m }
  {
    \bool_if:NTF \l_@@_thesis_in_declarations_bool
      { \@@_thesis_fm_part:nnn {#1} {#2} {#3} }
      { \msg_error:nnnn { tudcd } { thesis/fm-usage } { nextdeclaration } { declarations } }
  }
%</body&peri>
%    \end{macrocode}
